\section{Diagnóstico Institucional e Regulatório}

Este módulo concentra-se no mapeamento sistemático do ambiente institucional e regulatório de proteção aos conhecimentos tradicionais, com ênfase na identificação de gargalos de governança e custos de transação que impedem o acesso efetivo das comunidades aos marcos legais vigentes. A abordagem ancora-se na Economia dos Custos de Transação \cite{Williamson1985}, que concebe fricções institucionais como barreiras tanto ex ante quanto ex post. As barreiras ex ante manifestam-se na busca de informação, negociação e redação de contratos, ao passo que as barreiras ex post emergem no monitoramento, enforcement e resolução de disputas. Em contextos de conhecimentos tradicionais, tais fricções materializam-se na complexidade dos requisitos para registro de marcas coletivas, indicações geográficas e proteções sui generis, frequentemente agravadas pela baixa capacidade institucional das comunidades.

\subsection{Delineamento e Procedimentos}

O diagnóstico será conduzido mediante análise documental abrangente dos marcos legais aplicáveis, incluindo a Lei de Propriedade Industrial, o Protocolo de Nagoia e a legislação estadual pertinente, complementada por revisão detalhada dos normativos do INPI e ICMBio e dos requisitos para obtenção de certificações orgânicas, Fair Trade e Rainforest Alliance. Nessa etapa serão mapeados os documentos exigidos, prazos regulamentares, custos diretos tais como taxas e emolumentos, custos indiretos associados à assessoria técnica, além dos critérios de elegibilidade aplicáveis.

Concomitantemente, proceder-se-á à construção de diagramas de fluxo representando cada etapa dos processos de registro e certificação, com identificação dos pontos de decisão críticos, exigências documentais específicas e interfaces com órgãos externos. Técnicas de modelagem de processos de negócio serão empregadas para representação visual, permitindo rastreabilidade e transparência procedimental.

Entrevistas semiestruturadas serão conduzidas com gestores públicos vinculados ao INPI, SEBRAE e secretarias estaduais, bem como com assessores técnicos atuantes junto a comunidades quilombolas e advogados especializados em propriedade intelectual coletiva, totalizando entre 10 e 15 participantes. As entrevistas explorarão percepções sobre barreiras procedimentais, capacidade de resposta das instituições públicas e estratégias de mitigação adotadas por comunidades que alcançaram êxito em processos similares.

A estimação de custos de transação envolverá quantificação tanto de custos diretos quanto indiretos. Os custos diretos abrangem taxas oficiais, emolumentos cartorários e laudos técnicos obrigatórios, ao passo que os custos indiretos englobam horas de trabalho comunitário investidas, deslocamentos necessários e perda de renda durante trâmites burocráticos. Serão realizadas comparações entre custos enfrentados por comunidades formalmente organizadas e aquelas sem estrutura associativa consolidada, mensurando-se adicionalmente os tempos médios para conclusão de processos em diferentes cenários organizacionais.

\subsection{Análise de Dados e Produtos Esperados}

Os dados qualitativos oriundos de entrevistas e análise documental serão processados mediante análise temática \cite{Braun2006}, empregando codificação aberta e axial para identificação de categorias emergentes de barreiras, tais como déficits informacionais, acesso limitado a serviços especializados, morosidade burocrática crônica e assimetria de poder nas relações com órgãos reguladores. Os dados quantitativos referentes a custos e prazos serão organizados em tabelas comparativas e analisados mediante estatísticas descritivas, permitindo identificação de padrões e outliers.

Espera-se como resultado um mapeamento detalhado dos pontos críticos nos quais a intervenção do sistema proposto pode efetivamente reduzir custos de transação, acompanhado de um protocolo de governança simplificado capaz de sistematizar requisitos, organizar documentação necessária e orientar as comunidades no acesso estruturado a instrumentos de proteção jurídica de seus ativos intangíveis.
